% -*-coding: utf-8 -*-
%%%%%%%%%%%%%%%%%%%%%%%%%%%%%%%%%%%%%%%%%%%%%%%%%%%%%%%%%
\chapter{补充数学推导}[Supplementary Mathematical Derivations]

本附录补齐正文中压缩处理的推导细节，并给出与当前实现一致的矩阵表达。记号遵循正文：原子基序为 $(|0\rangle,|1\rangle,|e\rangle,|u\rangle)$，腔基序为 $(|{\rm vac}\rangle,|H\rangle,|V\rangle)$，bin 基序为 $(|{\rm vac}\rangle,|H_{780}\rangle,|V_{780}\rangle,|H_{1517}\rangle,|V_{1517}\rangle)$。

%%%%%%%%%%%%%%%%%%%%%%%%%%%%%%%%%%%%%%%%%%%%%%%%%%%%%%%%%
\section{从 QED 拉氏量到四能级有效哈密顿量}[From QED Lagrangian to the Four-Level Effective Hamiltonian]

\paragraph{近似链条与目标有效模型}

从
\begin{equation}
\mathcal{L}=\bar{\psi}(i\gamma^\mu\partial_\mu-m)\psi-eA_\mu\bar{\psi}\gamma^\mu\psi-\frac{1}{4}F_{\mu\nu}F^{\mu\nu}
\end{equation}
出发，在弱速、弱场和近共振条件下依次采用低能非相对论近似、原子内外自由度分离与绝热消元、长波偶极近似 $e^{i\mathbf{k}\cdot\mathbf{r}}\approx 1$、有限能级截断以及旋转波近似，可得有效哈密顿量
\begin{equation}
H_{\mathrm{eff}}(t)=H_{\mathrm{atom}}(t)+H_{\mathrm{cavity}}+H_{\mathrm{int}}+H_{\mathrm{diss}}.
\end{equation}

\paragraph{二能级极限与 Jaynes-Cummings 形式}

若先仅保留一条跃迁与一条腔模，则
\begin{equation}
H_{\mathrm{JC}}
=\frac{\omega_{ab}}{2}\sigma_z+\omega_c\left(a^\dagger a+\frac{1}{2}\right)
+g\left(\sigma^+a+\sigma^-a^\dagger\right).
\end{equation}
该式是四能级模型在“单激发、单偏振、单通道”极限下的退化情形，便于解释耦合常数 $g$ 与谱失谐参数的物理含义。

\paragraph{四能级原子矩阵与偏振跃迁算符}

驱动哈密顿量在基序 $(|0\rangle,|1\rangle,|e\rangle,|u\rangle)$ 下写为
\begin{equation}
H_{\mathrm{atom}}(t)=
\begin{bmatrix}
0&0&0&0\\
0&0&0&0\\
0&0&\Delta_e&\Omega(t)\\
0&0&\Omega^*(t)&\Delta_u
\end{bmatrix}.
\label{eq:app-hatom}
\end{equation}
偏振选择定则由
\begin{equation}
S_+=|0\rangle\langle e|=
\begin{bmatrix}
0&0&1&0\\
0&0&0&0\\
0&0&0&0\\
0&0&0&0
\end{bmatrix},
\quad
S_-=|1\rangle\langle e|=
\begin{bmatrix}
0&0&0&0\\
0&0&1&0\\
0&0&0&0\\
0&0&0&0
\end{bmatrix}
\end{equation}
给出。引入偏振映射矩阵
\begin{equation}
\boldsymbol{\alpha}=
\begin{bmatrix}
\alpha_{H,+} & \alpha_{H,-}\\
\alpha_{V,+} & \alpha_{V,-}
\end{bmatrix},
\end{equation}
则
\begin{equation}
\sigma_H=\alpha_{H,+}S_++\alpha_{H,-}S_-,
\qquad
\sigma_V=\alpha_{V,+}S_++\alpha_{V,-}S_-.
\end{equation}

\paragraph{三维腔模与 12 维发射体哈密顿量}

腔湮灭算符为
\begin{equation}
a_H=
\begin{bmatrix}
0&1&0\\
0&0&0\\
0&0&0
\end{bmatrix},
\qquad
a_V=
\begin{bmatrix}
0&0&1\\
0&0&0\\
0&0&0
\end{bmatrix}.
\end{equation}
于是
\begin{align}
H_{\mathrm{em}}(t)=&
H_{\mathrm{atom}}(t)\otimes I_3+I_4\otimes
\mathrm{diag}(0,\delta_{c,H},\delta_{c,V})\notag\\
&+g\left(
\sigma_H\otimes a_H^\dagger+\sigma_H^\dagger\otimes a_H
+\sigma_V\otimes a_V^\dagger+\sigma_V^\dagger\otimes a_V
\right).
\label{eq:app-hem}
\end{align}
可见输出算符为
\begin{equation}
L_{{\rm out},H}=\sqrt{\kappa_{{\rm ex},H}}\,I_4\otimes a_H,\quad
L_{{\rm out},V}=\sqrt{\kappa_{{\rm ex},V}}\,I_4\otimes a_V.
\end{equation}
不可见塌缩算符集合
\begin{equation}
\mathcal{C}=\left\{
\sqrt{\kappa_{{\rm in},H}}I_4\!\otimes\!a_H,\;
\sqrt{\kappa_{{\rm in},V}}I_4\!\otimes\!a_V,\;
\sqrt{\gamma_+}S_+\!\otimes\!I_3,\;
\sqrt{\gamma_-}S_-\!\otimes\!I_3
\right\}.
\end{equation}

%%%%%%%%%%%%%%%%%%%%%%%%%%%%%%%%%%%%%%%%%%%%%%%%%%%%%%%%%
\section{时间 bin 离散与发射门构造}[Time-Bin Discretization and Emission-Gate Construction]

\paragraph{连续场到离散场算符}

定义
\begin{equation}
b_{\mu,n}=\frac{1}{\sqrt{\Delta t}}\int_{t_n}^{t_n+\Delta t}b_\mu(t)\,dt,\qquad \mu\in\{H,V\},
\end{equation}
满足
\begin{equation}
[b_{\mu,m},b_{\nu,n}^\dagger]=\delta_{\mu\nu}\delta_{mn}.
\end{equation}
在短时展开下，单步幺正近似为
\begin{equation}
U_n=\exp\!\left[
\sqrt{\Delta t}\sum_{\mu}
\left(L_{{\rm out},\mu}\otimes b_{\mu,n}^\dagger-L_{{\rm out},\mu}^\dagger\otimes b_{\mu,n}\right)
-i\Delta t\,H_{\mathrm{em}}(t_n)\otimes I
\right].
\end{equation}

\paragraph{36 维子门嵌入 60 维 bin 门}

发射门先在 $(12\mathrm{D\ emitter}\otimes3\mathrm{D}_{780})$ 构造得到 $U_{36}$，再嵌入 $(12\mathrm{D}\otimes5\mathrm{D\ bin})$：
\begin{equation}
U_{60}=
\begin{bmatrix}
U_{36} & 0\\
0 & I_{24}
\end{bmatrix},
\end{equation}
其中 $I_{24}$ 对应 1517 子块保持单位。该结构保证“发射”不直接生成 1517 光子；1517 分量仅由后续 QFC 产生。

\paragraph{Kraus 轨迹更新与局域截断}

对不可见塌缩集合 $\mathcal{C}$，每步可写
\begin{equation}
\rho\mapsto\sum_j K_j\rho K_j^\dagger,\qquad \sum_j K_j^\dagger K_j=I.
\end{equation}
实现时采用单轨迹抽样并在 MPS 局域重分解，以维持计算复杂度可控。与全密度矩阵演化相比，该方案在保持物理口径一致的条件下显著降低内存与时间复杂度。

%%%%%%%%%%%%%%%%%%%%%%%%%%%%%%%%%%%%%%%%%%%%%%%%%%%%%%%%%
\section{QFC 与滤波记忆模的张量表达}[Tensor Representation of QFC and Filter Memory Modes]

\paragraph{QFC 的 5 维矩阵形式}

QFC 在 bin 基序下的显式矩阵为
\begin{equation}
U_{\mathrm{QFC}}=
\begin{bmatrix}
1 & 0 & 0 & 0 & 0\\
0 & c_H & 0 & -e^{-i\phi_H}s_H & 0\\
0 & 0 & c_V & 0 & -e^{-i\phi_V}s_V\\
0 & e^{i\phi_H}s_H & 0 & c_H & 0\\
0 & 0 & e^{i\phi_V}s_V & 0 & c_V
\end{bmatrix},
\end{equation}
其中 $c_\mu=\cos\theta_\mu,\ s_\mu=\sin\theta_\mu,\ \eta_\mu=\sin^2\theta_\mu$。

\paragraph{滤波腔离散系数与步进幺正}

一阶滤波腔离散参数
\begin{equation}
r=\exp\!\left[-(\pi\Delta\nu+i2\pi\delta)\Delta t\right],\qquad
t=\sqrt{1-|r|^2}
\end{equation}
满足能量守恒约束 $|r|^2+t^2=1$。

在 $(\mathrm{bin5D}\otimes\mathrm{mem3D})$ 的 $15$ 维空间中，除两个偏振子块外其余分量均为单位。定义
\begin{equation}
|h_1\rangle=|H_{1517},{\rm vac}_{\rm mem}\rangle,\quad
|h_2\rangle=|{\rm vac}_{\rm bin},H_{\rm mem}\rangle,
\end{equation}
\begin{equation}
|v_1\rangle=|V_{1517},{\rm vac}_{\rm mem}\rangle,\quad
|v_2\rangle=|{\rm vac}_{\rm bin},V_{\rm mem}\rangle.
\end{equation}
则
\begin{equation}
U_{\mathrm{step}}|_{\{h_1,h_2\}}=
\begin{bmatrix}
r&-t\\
t&r^*
\end{bmatrix},
\qquad
U_{\mathrm{step}}|_{\{v_1,v_2\}}=
\begin{bmatrix}
r&-t\\
t&r^*
\end{bmatrix}.
\label{eq:app-filter-sub}
\end{equation}

\paragraph{记忆链与最终布局}

每臂追加一个 mem 站点，扫描所有 bin 后收拢到链首得到固定布局
\begin{equation}
{\rm atomA,\ atomB,\ memA,\ memB,\ A1,\ B1,\ldots, AN,\ BN}.
\end{equation}
后续推导均基于该显式记忆路径，不再引入并行链路表示。

%%%%%%%%%%%%%%%%%%%%%%%%%%%%%%%%%%%%%%%%%%%%%%%%%%%%%%%%%
\section{Heisenberg-POVM 推回与统计量推导}[Heisenberg-POVM Pushback and Metric Derivations]

\paragraph{对偶映射恒等式}

对任意信道 $\Lambda(\rho)=\sum_kK_k\rho K_k^\dagger$ 与 effect $E$，
\begin{align}
\mathrm{Tr}\!\left[\Lambda(\rho)E\right]
&=\sum_k\mathrm{Tr}\!\left(K_k\rho K_k^\dagger E\right)\notag\\
&=\sum_k\mathrm{Tr}\!\left(\rho K_k^\dagger E K_k\right)\notag\\
&=\mathrm{Tr}\!\left[\rho\,\Lambda^\dagger(E)\right],
\end{align}
其中
\begin{equation}
\Lambda^\dagger(E)=\sum_k K_k^\dagger E K_k.
\label{eq:app-adjoint}
\end{equation}
因此可把光纤损耗、偏振旋转、相位噪声、探测效率与暗计数都推回到 effect，不必在态端引入混态膨胀。

\paragraph{可分辨性混合与背景 OR 映射}

记不可分辨与可分辨 effect 分别为 $E_{\mathrm{int}},E_{\mathrm{dist}}$，则
\begin{equation}
E=v_{\mathrm{res}}E_{\mathrm{int}}+(1-v_{\mathrm{res}})E_{\mathrm{dist}}.
\end{equation}
背景点击作为独立“OR”过程处理：若 $A$ 为信号点击集合、$B$ 为背景点击集合，观测集合 $S=A\vee B$，则
\begin{equation}
E_{\mathrm{obs}}[S]=\sum_A P(S|A)\,E_{\mathrm{sig}}[A].
\end{equation}
该式即“信号点击与背景点击独立 OR 合并”的一般表达。

\paragraph{成功概率、条件态与 Bell 保真度}

设成功记录集合为 $\mathcal{S}$，则
\begin{equation}
p_s=\sum_{m\in\mathcal{S}}\mathrm{Tr}(E_m\rho).
\end{equation}
条件态写为
\begin{equation}
\rho_{\mathrm{cond}}=
\frac{1}{p_s}
\sum_{m\in\mathcal{S}} M_m\rho M_m^\dagger.
\end{equation}
对目标 Bell 态 $|\Psi_{\rm tgt}\rangle$，条件保真度
\begin{equation}
F_{\mathrm{cond}}=\langle\Psi_{\rm tgt}|\rho_{\mathrm{cond}}|\Psi_{\rm tgt}\rangle.
\end{equation}
同时报告未归一化保真度
\begin{equation}
F_{\mathrm{full}}=\mathrm{Tr}\!\left(\rho_{\mathrm{cond}}\right)F_{\mathrm{cond}},
\end{equation}
用于反映成功概率与态质量的联合权重。

\paragraph{CHSH 最大值计算}

对两比特态 $\rho$，相关矩阵
\begin{equation}
T_{ij}=\mathrm{Tr}\!\left[\rho\left(\sigma_i\otimes\sigma_j\right)\right],\qquad i,j\in\{x,y,z\}.
\end{equation}
令 $U=T^{\mathsf T}T$，其最大两本征值为 $u_1,u_2$，则
\begin{equation}
S_{\max}=2\sqrt{u_1+u_2}.
\end{equation}
该式用于由收缩结果直接给出 CHSH 上界，不依赖额外抽样重构。

%%%%%%%%%%%%%%%%%%%%%%%%%%%%%%%%%%%%%%%%%%%%%%%%%%%%%%%%%
\section{单位口径与离散化一致性}[Unit Conventions and Discretization Consistency]

\paragraph{统一口径}

为避免“数值正确但物理时标错误”，统一规定：
\begin{equation}
\Omega,\ g,\ \Delta \ \text{采用}\ \mathrm{rad/s},\qquad
\kappa,\ \gamma \ \text{采用}\ \mathrm{1/s}.
\end{equation}
所有参数在进入发射内核前一次性换算，不允许在子模块重复隐式变换。

\paragraph{$\Delta t$ 变化下的不变量}

当离散步长 $\Delta t$ 调整时，连续模型中与时间积分相关的无量纲组合应保持一致，例如
\begin{equation}
g\sqrt{\Delta t},\quad \kappa\Delta t,\quad \gamma\Delta t,
\end{equation}
从而确保“改 dt 不改物理”的数值一致性。此原则是后续参数扫描与跨配置比较的基础。
